
\documentclass{article}
\usepackage{colortbl}
\usepackage{makecell}
\usepackage{multirow}
\usepackage{supertabular}

\begin{document}

\newcounter{utterance}

\centering \large Interaction Transcript for game `bbh', experiment `formal\_fallacies', episode 143 with Qwen3.6{-}27B.
\vspace{24pt}

{ \footnotesize  \setcounter{utterance}{1}
\setlength{\tabcolsep}{0pt}
\begin{supertabular}{c@{$\;$}|p{.15\linewidth}@{}p{.15\linewidth}p{.15\linewidth}p{.15\linewidth}p{.15\linewidth}p{.15\linewidth}}
    \# & \multicolumn{2}{c}{Player} && \multicolumn{2}{c}{Game Master} \\
    \hline

    \theutterance \stepcounter{utterance}  
    & & & \multicolumn{4}{p{0.6\linewidth}}{
        \cellcolor[rgb]{0.9,0.9,0.9}{
            \makecell[{{p{\linewidth}}}]{
                \texttt{\tiny{[P1$\langle$GM]}}
                \texttt{Distinguish deductively valid arguments from formal fallacies.} \\
\\ 
\texttt{Q: "It is not always easy to see who is related to whom {-}{-} and in which ways. The following argument pertains to this question: To begin with, Lesley is a close friend of Fernando. Moreover, being a close friend of Fernando or a schoolmate of Lowell is sufficient for being a great{-}grandfather of Leroy. It follows that Lesley is a great{-}grandfather of Leroy."} \\
\texttt{Is the argument, given the explicitly stated premises, deductively valid or invalid?} \\
\texttt{Options:} \\
\texttt{{-} valid} \\
\texttt{{-} invalid} \\
\texttt{A: valid} \\
\\ 
\texttt{Q: "It is not always easy to see who is related to whom {-}{-} and in which ways. The following argument pertains to this question: Whoever is not a great{-}grandfather of Clyde is a stepbrother of Brian. Being an ancestor of Dana is sufficient for not being a great{-}grandfather of Clyde. We may conclude: Everyone who is an ancestor of Dana is a stepbrother of Brian, too."} \\
\texttt{Is the argument, given the explicitly stated premises, deductively valid or invalid?} \\
\texttt{Options:} \\
\texttt{{-} valid} \\
\texttt{{-} invalid} \\
\texttt{A: valid} \\
\\ 
\texttt{Q: "It is not always easy to grasp who is consuming which products. The following argument pertains to this question: Every infrequent user of Paul Mitchell shampoo is either a rare consumer of Nioxin shampoo or a loyal buyer of Caress soap, or both. No regular consumer of Lush soap is a rare consumer of Nioxin shampoo and, in the same time, a loyal buyer of Caress soap. It follows that whoever is an infrequent user of Paul Mitchell shampoo is not a regular consumer of Lush soap."} \\
\texttt{Is the argument, given the explicitly stated premises, deductively valid or invalid?} \\
\texttt{Options:} \\
\texttt{{-} valid} \\
\texttt{{-} invalid} \\
\texttt{A: invalid} \\
\\ 
\texttt{Q: "Is Siri a stepsister of Mary? Is Susan related to Kate? In large families, it is sometimes difficult to keep track of all one's relatives. The following argument seeks to clarify some such relations: To start with, Merle is a cousin of Annie. Now, whoever is not both a classmate of Hazel and a cousin of Annie is a great{-}grandmother of Doria. We may conclude that it is false that Merle is a great{-}grandmother of Doria."} \\
\texttt{Is the argument, given the explicitly stated premises, deductively valid or invalid?} \\
\texttt{Options:} \\
\texttt{{-} valid } \\
\texttt{{-} invalid} \\
\texttt{A:} \\
            }
        }
    }
     \\ \\

    \theutterance \stepcounter{utterance}  
    & \multicolumn{4}{p{0.6\linewidth}}{
        \cellcolor[rgb]{0.561,0.737,0.561}{
            \makecell[{{p{\linewidth}}}]{
                \texttt{\tiny{[P1$\rangle$GM]}}
                \texttt{To determine the validity of the argument, let's break it down into logical components.} \\
\\ 
\texttt{**1. Identify the Premises and Conclusion:**} \\
\\ 
\texttt{*   **Premise 1:** Merle is a cousin of Annie.} \\
\texttt{    *   Let \$M\$ = Merle} \\
\texttt{    *   Let \$C\_A\$ = is a cousin of Annie} \\
\texttt{    *   Statement: \$C\_A(M)\$ is True.} \\
\\ 
\texttt{*   **Premise 2:** Whoever is not both a classmate of Hazel and a cousin of Annie is a great{-}grandmother of Doria.} \\
\texttt{    *   Let \$C\_H\$ = is a classmate of Hazel} \\
\texttt{    *   Let \$G\_D\$ = is a great{-}grandmother of Doria} \\
\texttt{    *   Logical Form: For any person \$x\$, if \$\textbackslash{}neg (C\_H(x) \textbackslash{}land C\_A(x))\$, then \$G\_D(x)\$.} \\
\texttt{    *   Symbolically: \$\textbackslash{}forall x, \textbackslash{}neg (C\_H(x) \textbackslash{}land C\_A(x)) \textbackslash{}rightarrow G\_D(x)\$} \\
\\ 
\texttt{*   **Conclusion:** It is false that Merle is a great{-}grandmother of Doria.} \\
\texttt{    *   Statement: \$\textbackslash{}neg G\_D(M)\$} \\
\\ 
\texttt{**2. Analyze the Logic:**} \\
\\ 
\texttt{We need to determine if the conclusion \$\textbackslash{}neg G\_D(M)\$ necessarily follows from the premises.} \\
\\ 
\texttt{From Premise 1, we know \$C\_A(M)\$ is True.} \\
\\ 
\texttt{Let's look at} \\
            }
        }
    }
    & & \\ \\

    \theutterance \stepcounter{utterance}  
    & & & \multicolumn{2}{p{0.3\linewidth}}{
        \cellcolor[rgb]{0.9,0.9,0.9}{
            \makecell[{{p{\linewidth}}}]{
                \texttt{\tiny{[GM$|$GM]}}
                \texttt{Target: invalid} \\
            }
        }
    }
    & & \\ \\

    \theutterance \stepcounter{utterance}  
    & & & \multicolumn{2}{p{0.3\linewidth}}{
        \cellcolor[rgb]{0.9,0.9,0.9}{
            \makecell[{{p{\linewidth}}}]{
                \texttt{\tiny{[GM$|$GM]}}
                \texttt{game\_result = LOSE} \\
            }
        }
    }
    & & \\ \\

\end{supertabular}
}

\end{document}
